\section{Conclusion}
\label{sec:conclusion}

The matched comparison shows that the relative performance of TimesNet and
DilatedRNN depends on the GHI forecasting task. Across ten factory-associated
locations, inputs available at the forecast origin, data partitions,
forecast origins, post-processing, metrics, and location weights were
controlled within each task. The results describe forecasts of the NSRDB solar resource, without
quantifying factory demand or photovoltaic output.

For cumulative macro-MAE, TimesNet led at 24 and 72 hours and at the daily
checkpoints up to 14 days. DilatedRNN led at 48 hours, 30 days, and in the
primary one-month task. In the exploratory six-month task, TimesNet led at
the one-month prefix, while DilatedRNN led at the three- and six-month
prefixes. Local win counts sometimes disagreed with the macro ranking:
DilatedRNN achieved lower terminal MAE at six of the ten sites in each of
the hourly, daily, and primary one-month tasks, and at seven sites in the
exploratory six-month task.

The complete benchmark puts these architecture-specific advantages in
perspective. XGBoost led all hourly checkpoints and the daily checkpoints
at 1, 3, and 14 days; LSTM led at seven days, and DilatedRNN led at 30 days
with 43.85~W/m$^2$, compared with 43.93~W/m$^2$ for XGBoost. Calendar
climatology led the monthly checkpoints. Neither TimesNet nor DilatedRNN
was consistently best across tasks.

Daily and monthly estimates use five-seed ensembles, and several differences
are small relative to the observed seed variability. The single hourly seed,
overlapping targets, one test year, and unequal parameter budgets limit
generalization. Post-hoc weather cases provide context but do not explain
model differences causally.

Future evaluation should include prospective years and unseen sites,
multiple hourly seeds, daylight-only metrics, measured computational cost,
and weather forecasts available at issuance. The matched protocol and
origin-level predictions provide a basis for testing these extensions.
Together, the task-level rankings, site-level differences, and reference
benchmarks show why architecture choice must be evaluated for the intended
forecasting task.
